\chapter{Konzept und Architektur}\label{ch:concept}

Dieses Kapitel präsentiert den Kernbeitrag der Arbeit: ein Achsen-Bibliotheks-Framework, das
cache-bewusste Suchalgorithmen in austauschbare, orthogonale Bausteine-Achsen zerlegt, die
geschichtete Engine-Architektur, die es realisiert, die ABI-stabile Modul-Schnittstelle, die
neunzehn-Achsen-Permutationsmatrix, den Prüfling PRT-ART und den Builder, der die
Drei-Stufen-Prüfung orchestriert.

\section{Das Achsen-Bibliotheks-Framework}\label{sec:axis-framework}

Die zentrale Idee besteht darin, einen Algorithmus nicht als Monolith zu behandeln, sondern als
\emph{Komposition orthogonaler Achsen}, wobei jede Achse eine Teilentscheidung eines cache-bewussten
Index repräsentiert (wie Seiten angeordnet werden, wie Knoten verzweigen, wie Speicher alloziert
wird, wie Nebenläufigkeit behandelt wird usw.). Ein konkreter Algorithmus ist dann ein Punkt im
kartesischen Produkt aller Achsen, und ein Stand-der-Technik-Entwurf wird als eine explizite
Konfiguration rekonstruiert. Die entscheidende Eigenschaft ist die \emph{modulare Austauschbarkeit}:
Forschende können eine neue Achsen-Variante isoliert untersuchen --- und ihren Effekt pro Achse und
für den gesamten Algorithmus messen --- ohne alle anderen neu zu implementieren. Das macht den
Vergleich zweier \texttt{std::map}-Implementierungen (Abschnitt~\ref{sec:stdmap-interface}) zu einem
kontrollierten, reproduzierbaren Experiment.

Diese Zerlegung beruht auf zwei Beobachtungen. Erstens teilen Suchalgorithmen im großen Bild
dieselbe \emph{Baum-Anatomie} in unterschiedlicher Ausprägung --- von ausgewogenen Suchbäumen über
Tries und Radix-Bäume bis zu eigenentwickelten Konstrukten; selbst eine sortierte
Schlüssel-Wert-Liste ist streng genommen ein einzelnes Blatt eines Baumes, dessen Werte über
Verweise weiter auf Daten-Blätter zeigen. Die Baum-Struktur ist damit universell über alle
Suchalgorithmus-Architekturen hinweg und bildet die Leinwand, auf der die Achsen platziert werden.
Nicht-baumartige Suchverfahren --- etwa eine Hash-Tabelle --- fügen sich in dasselbe Achsen-System
ein, indem sie die baum-spezifischen Achsen (Pfadkompression, Knotentyp, Bereichs-Index) als
Durchreich-Varianten belegen und nur das gemeinsame Achsen-Subset tragen; sie bleiben dieselbe
Gattung mit eingeschränktem Profil. Eine eigene Gattung entstünde erst, wenn sich der
Basis-Achsen-Satz grundlegend unterschiede. Zweitens fasst eine anschauliche Metapher einen
konkreten Algorithmus als \emph{Lebewesen} auf, dessen \emph{Organe} die Achsen sind: Jeder
Algorithmus besitzt ein Knoten-Layout-Organ, ein Traversal-Organ, ein Allokations-Organ und so
weiter, jeweils in anderer Ausprägung. Die Permutation der Achsen gleicht dann einem genetischen
Experiment --- der Austausch von Organen zwischen Lebewesen erzeugt neue Komposit-Algorithmen mit
bisher unbekanntem Cache-Line-Verhalten. Das Forschungsziel ist es, diese gemeinsame Anatomie und im
resultierenden Entwurfsraum
(\emph{design space})~\cite{idreos2018periodic, idreos2018datacalculator} das schnellste Komposit
für ein gegebenes datentyp-spezifisches Lastmuster zu finden.

\section{Vier-Subsystem-Trennung (M-Modell)}\label{sec:m-model}

Die Diplomarbeit, der Cache-Engine-Builder, die Cache-Engine und der Prüfling sind
\emph{vier formal getrennte Subsysteme} mit unterschiedlichen Verantwortlichkeiten. Die Trennung ist
orthogonal zur Drei-Repository-Aufteilung (Abschnitt~\ref{sec:repos}): Builder und Engine wohnen im
selben Repository, sind dort aber eine ausführbare Datei bzw.\ eine Bibliothek.

\begin{enumerate}
  \item \textbf{\texttt{messung\_driver}} (Diplomarbeit/Code) --- der Auswertungs-Orchestrator
  (äußere Schleife): liest die drei Messreihen-XML-Konfigurationen (A/B/C), triggert pro Reihe den
  Builder und sammelt die Ergebnisse für den Anhang.
  \item \textbf{CacheEngineBuilder} (\texttt{cache-engine/apps}) --- ein autonomes
  Plattform-Ausmess-System, das die Sieben-Phasen-Pipeline implementiert (discover, measure,
  classify, publish, bind, execute, compare) und je registrierter Engine den Permutationsraum
  enumeriert.
  \item \textbf{CacheEngine} (\texttt{cache-engine/libs}) --- die Werkzeug-Bibliothek: sie stellt
  für jede der neunzehn Bausteine-Achsen (Abschnitt~\ref{sec:axes}) die Standard-Implementierungen
  und die zugehörigen Cache-Strategie-Familien als C++23-Concept-API bereit, aus denen der Builder
  die Permutations-Konfigurationen zusammensetzt.
  \item \textbf{Prüfling (PRT-ART)} (\texttt{prt-art}) --- eine bei der Cache-Engine als
  \texttt{IExecutionEngine} registrierte Algorithmus-Implementierung; sie stellt ihre eigenen
  Achsen-Organe bereit und konsumiert Cache-Engine-Dienste während ihrer eigenen
  Lookup-/Insert-/Scan-Operationen.
\end{enumerate}

Die definierende Eigenschaft des M-Modells ist die \emph{bidirektionale}
Cache-Engine~$\leftrightarrow$~Prüfling-Beziehung: (a) die Cache-Engine instantiiert den Prüfling je
Permutations-Konfiguration (Builder-Phase 5, \textsc{bind}), und (b) der Prüfling ruft zur Laufzeit
Cache-Engine-Dienste (Telemetry, Prefetch, Heuristik) während seiner eigenen Lookups auf (Phase 6,
\textsc{execute}). Diese Symmetrie ermöglicht die \emph{Multi-Prüfling}-Fähigkeit: Messreihe~A
vergleicht PRT-ART gegen viele SOTA-Adapter in einer einzigen Builder-Pipeline.

\section{Die eine Architektur}\label{sec:three-layer}

Das Modell kennt \emph{eine} Architektur --- eine einzige, in sich geschlossene Hierarchie ohne
Parallel-Bäume. An ihrer Wurzel steht die \texttt{ExecutionEngine} als gemeinsame, ausmessbare
Abstraktion über allem Messbaren; unter ihr verzweigen sich lediglich \emph{Lebewesen} (Anatomien
mit Achsen-Organen) und achsenlose \emph{Viren} (etwa reine Graphen-Algorithmen) als Geschwister
derselben Wurzel.

\footnotesize
\begin{verbatim}
    IExecutionEngine                               Wurzel: alles Ausmessbare
      |
      +-- IAnatomyBase : IExecutionEngine          Lebewesen
      |     |
      |     +-- SearchAlgorithm-Unterklasse        = "Lebewesen" (Saeugetier)
      |           |
      |           +-- SearchAlgorithmAnatomy<C>     = sein KOERPER (19 Organe)
      |           +-- SearchAlgorithmAbiAdapter<A>  = ABI-Laufzeit-SICHT ("SearchEngine")
      |
      +-- IVirusExecutionEngine : IExecutionEngine  Virus (achsenlos, Graph)
\end{verbatim}
  \normalsize

  Ein Suchalgorithmus ist in diesem Modell ein \emph{Lebewesen}, und „Lebewesen" und
  \texttt{SearchAlgorithm} bezeichnen denselben Gegenstand --- Metapher und technischer Begriff sind
  zwei Namen für eine Sache. Nach außen wird das Lebewesen über sein Gattungs-Interface (das
  Prüf-Dock, ein \texttt{std::map}-artiges Interface) angesprochen; sein Inneres ist seine
  \emph{Anatomie}. Die Anatomie ist nicht ein zweites Modell neben dem Suchalgorithmus, sondern sein
  \emph{Körper}: die feste Komposition seiner neunzehn Achsen, von denen jede einem Organ
  entspricht. Über die Modulgrenze hinweg wird dasselbe Lebewesen als ABI-Sicht sichtbar gemacht ---
  die ABI-/Laufzeit-Schichtung
  \texttt{CacheEngine}\,$\to$\,\texttt{ExecutionEngine}\,$\to$\,\texttt{SearchEngine} ist genau
  diese Sicht desselben Lebewesens über die Modulgrenze, kein paralleles Konstrukt:
  \texttt{CacheEngine} stellt die Standard-Bausteine der Achsen als Bibliothek,
  \texttt{ExecutionEngine} die OS-Primitiven (cache-line-aligned Allokation, NUMA-konformes
  Read-Pin, coherence-aware Write), und „SearchEngine" ist die Such-API-Sicht desselben Lebewesens,
  gemessen über genau eine Beobachter-Schnittstelle über alle neunzehn Organe. Es gibt somit eine
  einzige Kette --- Wurzel (\texttt{ExecutionEngine}) über das Gattungs-Prüf-Dock zu dem Lebewesen,
  das gleichbedeutend mit \texttt{SearchAlgorithm} ist und dessen Körper die Anatomie mit neunzehn
  Achsen ist; jede scheinbare Dopplung in der Implementierung ist eine zu vereinheitlichende Stelle,
  kein gewolltes Nebeneinander.

  \section{ABI-stabiles C++23-Modul-Interface}\label{sec:abi}

  Das ABI folgt dem Prinzip variadischer Template-Parameter: ein Parameter $\Rightarrow$
  \texttt{std::vector}-API, zwei Parameter $\Rightarrow$ \texttt{std::map}-API, $N>2 \Rightarrow$
  \texttt{std::map<K,\,std::tuple<V$_1$\ldots V$_N$>>}. Komplexe Schlüssel werden von einer
  Fingerprint-Komponente auf 16-Byte-Binärstrings fester Länge reduziert. Die Modulgrenze ist
  binär-stabil, sodass der Builder vorkompilierte Permutations-Module laden (LoadLibrary/dlopen) und
  den ABI-Vertrag je Modul verifizieren kann.

  \section{Die Bausteine-Achsen}\label{sec:axes}

  Nach der N-Phasen-Erweiterung und der Konsolidierung zur Suchalgorithmus-Anatomie umfasst die
  Matrix \textbf{neunzehn Hauptachsen plus etwa siebenundfünfzig Sub-Achsen}: (1)~Search-Algorithm,
  (2)~Cache-Traversal, (3)~Mapping, (4)~Path-Compression, (5)~Node-Type, (6)~Memory-Layout,
  (7)~Allocator (Allocation, Reclamation, NUMA, Huge-Page, Free-List), (8)~Prefetch, (9)~Concurrency
  (Pattern, Locking-Mode), (10)~Serialization, (11)~Telemetry, (12)~Value-Handle, (13)~ISA,
  (14)~Index-Organization, (15)~I/O-Dispatch, (16)~Migration, (17)~Filter, (18)~Queuing-q1 und
  (19)~Queuing-q2 --- die siebzehn Such-Achsen plus die beiden Queuing-Achsen, genau die neunzehn
  Organe, die das \texttt{IsComposition}-Concept je Algorithmus erzwingt. Je Achse enumeriert ein
  \texttt{std::variant} alle Konkretisierungen (allein die Allocator-Achse bietet rund zwei
  Dutzend); das kartesische Produkt ergibt weit mehr als $10^{11}$ Konfigurationen
  (Abschnitt~\ref{sec:explosion}).

  \section{PRT-ART als Prüfling}\label{sec:prt-art}

  PRT-ART ist ein \emph{Prüfling}: ein abstraktes Lebewesen, das einige Achsen-Organe selbst stellt
  und für die übrigen über einen Compile-Time-Fallback (\texttt{resolve\_baustein.hpp}) auf die
  Cache-Engine-Standardbausteine zurückgreift, wo es eine Achse nicht überschreibt; die hybride
  \texttt{PrtArtSearchEngine<Ts\ldots>}-API bleibt davon unberührt. Die Laufzeit-Anbindung erfolgt
  über den Execution-Engine-Adapter, nicht über eine eigene Such-Engine-Schicht. PRT-ART steuert
  eigene Implementierungen in den Achsen Page, Layout, Free-List, Prefetch, Concurrency-Pattern und
  Measurement bei (4+2-Pool-Allokator, Distance-Estimator-Prefetch, OLC mit reservierten
  Wert-Blöcken sowie die H1/H2/H3-Metriken); für die übrigen Achsen verwendet er bestehende
  SOTA-Bausteine über die Permutations-Konfiguration wieder.

  \section{CacheEngineBuilder und Drei-Stufen-Prüfung}\label{sec:builder}

  Der Builder orchestriert eine Drei-Stufen-Prüfung: \textbf{Stufe~1} (nur Cache-Engine) nutzt die
  Standard-Variante je Achse; \textbf{Stufe~2} (nur Prüfling) ersetzt die überschriebenen Achsen
  durch die Prüfling-Varianten und behält sonst die Cache-Engine-Standards; \textbf{Stufe~3} (Full
  Join, nicht-redundant) kombiniert die Standard- und alle Prüfling-Varianten. Jede Stufe erzeugt
  tausende nicht-redundante Permutations-Binaries, die zur Laufzeit jeweils als ABI-stabiles
  C++23-Modul geladen werden. Die Stufen bilden direkt die drei Pflicht-Messreihen aus
  Kapitel~\ref{ch:methodology} ab. Aus den feingliedrigen Messergebnissen liefert die Pipeline
  bereits die nötige Rangbildung; die automatische Auswahl des besten Gesamtalgorithmus und seine
  Auslieferung als eigenständige, ABI-stabile Binary hinter dem einheitlichen
  Suchalgorithmus-Interface sind als Erweiterung vorgesehen (Zielsetzung, Ausblick).
